\section{Conclusion and Recommendations}

This paper proposed and implemented a four stage framework that segments taxpayers into peer groups, detects anomalous reporting within each group using Isolation Forest, tests the learnability of documented anomaly mechanisms with supervised classifiers, and explains segment membership and anomaly flags with SHAP. On the ATO 2022--23 sample file, segmentation on a nine-dimensional income-composition space produced eight well-separated, tax-interpretable peer groups with silhouette score 0.65, showing that the value of taxpayer segmentation depends on choosing a space that describes what kind of return is filed rather than how much is reported.

The segment based Isolation Forest changed more than half of a 2\% review list relative to a population-wide forest and removed the systematic over-flagging of rare taxpayer types. The evaluation did not support the assumption that segmentation by itself makes the detector more sensitive: income confounding within the large wage segment persisted, and the ablation showed that the feature space matters at least as much as the reference population. Isolation Forest remains an appropriate primary detector for this data, fast, scale-invariant and exactly explainable, but its flags are dominated by rarely reported items and it is weak on single-feature inconsistencies. SHAP made both models inspectable at the level of the return, with a surrogate bridging the clustering stage, and every explanation was verified against the model's own score. For audit-risk assessment the message is that a segment-conditional, explainable anomaly score is a credible prioritisation tool, provided its output is treated as a signal for human review, its budget is set as a policy choice, and its sensitivity is measured against documented mechanisms rather than assumed.

\subsection{Recommendations and Future Research}

Future research should pursue hybrid designs in which unsupervised scores become features or weak labels for supervised models; deep autoencoders and graph-based detection linking taxpayers to agents, employers, partnerships and trusts~\cite{akoglu2015graph}; temporal modelling across years so that transitions are distinguished from irregularities; anomaly detection in a scale-free behavioural feature space, which the ablation indicates is the more promising direction for peer-group designs; fairness evaluation and bias mitigation so that review burden does not fall disproportionately on particular occupations, regions or age groups~\cite{barocas2023fairness}; and human-in-the-loop studies in which auditor feedback on explained flags is fed back into the models.
